\documentclass[hyphens,aspectratio=169,dvipsnames]{beamer}
\usepackage{graphicx}
\usepackage{xcolor}
\usepackage{amssymb}
\usepackage{amsmath}
\usepackage{mathtools}
\usepackage{textcomp}
\usepackage{moresize}
\usepackage{framed}
\usepackage{minted}
\usepackage{relsize}
\usepackage{tikz}
\usepackage{comment}
\usetikzlibrary{shapes.geometric, arrows, positioning}
\usetheme{Berlin}
\usecolortheme[RGB={0,95,47}]{structure}
\beamertemplatenavigationsymbolsempty

\makeatletter
\AtBeginEnvironment{minted}{\dontdofcolorbox}
\def\dontdofcolorbox{\renewcommand\fcolorbox[4][]{##4}}
\makeatother

\newcommand{\textpf}[1]{\texttt{\color{black}\fcolorbox{lightgray}{lightgray}{#1}}}

\begin{document}

\begin{frame}[fragile]{CS 15}
    \begin{center} Intro to Systems Programming \end{center}
\end{frame}

\begin{frame}[fragile]{Expectations}
    \begin{itemize}
        \pause \item This is a \textbf{closed laptop} class.
        \pause \item There will be in-class activities, during which computer use is permitted.
        \pause \item Outside of these activities, \textbf{do not use your laptop}.
        \pause \item During an activity, it is expected that your laptop is used \textbf{only for the activity}.
        \pause \item If we spot you using your laptop for unrelated stuff, we may publicly ask you to stop.
        \pause \item Repeat offenders will be asked to leave the classroom.
    \end{itemize}
\end{frame}


\begin{frame}[fragile]{C}
    \pause C is known for being
    \begin{itemize}
        \pause \item fast, % if you know what you're doing
        \pause \item portable, % if you know what you're doing
        \pause \item non-portable, % if you use platform-specific features, which is easy to do (intentionally and unintentionally)
        \pause \item low-level \pause (useful for ``close-to-the-metal" tasks),
        \pause \item old, % around since the 70s
        \pause \item hard to use correctly.
    \end{itemize}
\end{frame}

\begin{frame}[fragile]{Machine Code}
    \begin{itemize}
        \pause \item A computer's native language is known as ``machine code."
        \pause \item At the very bottom, your CPU's primary job is to interpret machine code.
        \pause \item Your CPU doesn't know how to execute Python, or Java, or C directly. It can only execute machine code.
    \end{itemize}
\end{frame}

\begin{frame}[fragile]{Interpreters}
    \begin{itemize}
        \pause \item If your CPU can only execute machine code, how do Python programs run?
        \pause \item Your Python programs are interpreted at runtime by a native machine code program (the Python interpreter).
        \pause \item This incurs overhead at runtime.
    \end{itemize}
\end{frame}

\begin{frame}[fragile]{Compilers}
    \begin{itemize}
        \pause \item A compiler is a program that translates programs from one language to another.
        \pause \item C programs are not usually interpreted directly.
        \pause \item Instead, they are translated into machine code using a compiler, and that machine code is interpreted by the CPU.
    \end{itemize}
\end{frame}

\begin{frame}[fragile]{Building C Programs}
    There are 4 tools required to turn a C program into an executable:
    \begin{enumerate}
        \pause \item The C preprocessor, which evaluates special ``directives" in your \textpf{.c} file. producing a \textpf{.i} file.
        \pause \item The C compiler, which translates the \textpf{.i} file into an assembly (\textpf{.s}) file.
        \pause \item The assembler, which translates the assembly into machine code in a \textpf{.o} file.
        \pause \item The linker, which links together \textpf{.o} file(s) into a single executable or library.
    \end{enumerate}
\end{frame}

\begin{frame}[fragile]{The Compiler Driver}
    \begin{itemize}
        \pause \item A C compiler driver is a program that invokes the preprocessor, compiler, assembler, and linker for us, so that we can build C programs all in one step.
        \pause \item There are 2 compiler drivers installed on Plink: \textpf{clang} and \textpf{gcc}.
        \pause \item Most of the time, these can be used interchangeably. Try both!
        \pause \item Often, people refer to the compiler driver as a ``compiler." This is synecdoche; it's okay that it's not strictly correct.
    \end{itemize}
\end{frame}

\begin{frame}[fragile]{A Simple C Program}
    Please copy down this C program into a file named \textpf{simple.c} on Plink:
    \inputminted{c}{./simple.c}
\end{frame}

\begin{frame}[fragile]{Explanation of \textpf{simple.c}}
    \begin{itemize}
        \pause \item \textpf{int}
            \begin{itemize}
                \pause \item The return type of the \textpf{main} function.
            \end{itemize}
        \pause \item \textpf{main}
            \begin{itemize}
                \pause \item The name of the \textpf{main} function.
            \end{itemize}
        \pause \item \textpf{(void)}
            \begin{itemize}
                \pause \item Indicates that \textpf{main} takes no arguments.
                \pause \item You might expect that empty parentheses would indicate no arguments, but it actually indicates that the function accepts any number of arguments, and ignores them all.
            \end{itemize}
    \end{itemize}
\end{frame}

\begin{frame}[fragile]{Running the Compiler Driver}
    \begin{itemize}
        \pause \item To use the \textpf{gcc} compiler driver to build \textpf{simple.c}, run
        \begin{minted}{sh}
            gcc ./simple.c
        \end{minted}
        \pause \item You should now have an executable named \textpf{a.out}.
        \pause \item To execute \textpf{a.out}, run
        \begin{minted}{sh}
            ./a.out
        \end{minted}
        \pause \item To see the returned value from \textpf{main}, run
        \begin{verbatim}
            echo "$?"
        \end{verbatim} % making this a minted would be nice, but editor isn't highlighting it properly due to the $.
    \end{itemize}
\end{frame}

\begin{frame}[fragile]{Integer Types in C}
    There are 5 basic integer types in C:
    \begin{itemize}
        \pause \item \textpf{char} (smallest),
        \pause \item \textpf{short} (small),
        \pause \item \textpf{int} (medium),
        \pause \item \textpf{long} (large).
        \pause \item \textpf{long long} (largest).
    \end{itemize}
\end{frame}

\begin{frame}[fragile]{Integer Sizes}
    \begin{itemize}
        \pause \item C is designed for portability, so the language standard doesn't specify the number of bytes required to store each of the integer types.
        \pause \item For this reason, we need the \textpf{sizeof} operator.
        \pause \item The \textpf{sizeof} operator returns the number of bytes used to represent its operand.
    \end{itemize}
\end{frame}

\begin{frame}[fragile]{Activity \#1: \textpf{sizeof}}
    Modify \textpf{simple.c} to determine the sizes of the 5 basic integer types on Plink.
\end{frame}

\begin{frame}[fragile]{Signedness}
    \begin{itemize}
        \pause \item All integers in C are either \textpf{signed} or \textpf{unsigned}.
        \pause \item Unsigned integers can represent only nonnegative values.
        \pause \item Signed integers can represent both negative and nonnegative values.
    \end{itemize}
\end{frame}

\begin{frame}[fragile]{Why would we ever want to use unsigned integers?}
    \begin{itemize}
        \pause \item They might better match the problem domain.
            \begin{itemize}
                \pause \item e.g., quantities, timestamps, etc.
            \end{itemize}
        \pause \item The positive range is larger!
            \begin{itemize}
                \pause \item The largest value that an \textpf{unsigned int} can store is greater than the largest value that a \textpf{signed int} can store.
            \end{itemize}
        \pause \item They are defined to wrap around.
            \begin{itemize}
                \pause \item Adding 1 to the biggest \textpf{unsigned int} is guaranteed to be \textpf{0}.
                \pause \item Adding 1 to the biggest \textpf{int} is \textbf{undefined}.
            \end{itemize}
        \pause \item Some style guides (e.g., Google's) say that you should never use unsigned types at all. Ask me about this if you're curious why.
    \end{itemize}
\end{frame}

\begin{frame}[fragile]{Integer Ranges}
    \begin{itemize}
        \pause \item 1 bit can store one of $2^1 = 2$ distinct values: \textpf{0} or \textpf{1}.
        \pause \item 2 bits can store one of $2^2 = 4$ distinct bitstrings: \textpf{00}, \textpf{01}, \textpf{10}, or \textpf{11}.
        \pause \item Therefore, \textpf{char}, which is a 1-byte (8-bit) integer type, can store one of $2^8 = 256$ distinct bitstrings.
        \pause \item For \textpf{unsigned char}, we treat those bitstrings as representing $0$-$255$, inclusive.
        \pause \item For \textpf{signed char}, we treat those bitstrings as representing $-128$-$127$, inclusive.
    \end{itemize}
\end{frame}

\begin{frame}[fragile]{Integer Literals}
    In C, integer literals (e.g., the \textpf{5} in \textpf{return x + 5;}) have the type \textpf{int}.

    \pause To construct literals of other integer types, append the following suffixes to the literal:
    \begin{itemize}
        \pause \item \textpf{l}: long
        \pause \item \textpf{ll}: long long
        \pause \item \textpf{u}: unsigned int
        \pause \item \textpf{ul}: unsigned long
        \pause \item \textpf{ull}: unsigned long long
    \end{itemize}

\end{frame}

\begin{frame}[fragile]{Integer Operators}
    \begin{itemize}
        \pause \item \textpf{+}: Addition
        \pause \item \textpf{-}: Subtraction
        \pause \item \textpf{*}: Multiplication
        \pause \item \textpf{/}: Division (rounds toward 0)
        \pause \item \textpf{\%}: Remainder
        \pause \item ...and there are more, but we'll cover them later :)
    \end{itemize}
\end{frame}

\begin{frame}[fragile]{Activity \#2: The math is all weird!}
    \begin{itemize}
        \pause \item These operators do not behave the way you expect!
        \pause \item Enumerate as many weird properties of C's math as you can.
        % Give them a few minutes, then collect these things on the board.
    \end{itemize}

\end{frame}

\begin{frame}[fragile]{Activity \#3: \textpf{abs}}
    \begin{center} Write a C function, \textpf{my\_abs(signed char x)}, that returns the absolute value of \textpf{x}. \end{center}

    \begin{itemize}
        \pause \item What is the most appropriate return type for your function?
            \begin{itemize}
                \pause \item \textpf{char}, \textpf{short}, \textpf{int}, or \textpf{long}?
                \pause \item \textpf{signed} or \textpf{unsigned}?
            \end{itemize}
        \pause \item Test your function on a wide range of inputs. Does it work? % If they chose (signed) char as their return type, they won't be able to run abs(-128).
    \end{itemize}
\end{frame}

\begin{frame}[fragile]{Activity \#4: $\log_2$}

    \begin{itemize}
        \pause \item Write a C function, \textpf{my\_log2(unsigned long x)}, that returns the $\log_2$ of \textpf{x}, rounded down to the nearest integer.
        \pause \item What should your function's return type be?
            \begin{itemize}
                \pause \item \textpf{char}, \textpf{short}, \textpf{int}, or \textpf{long}?
                \pause \item \textpf{signed} or \textpf{unsigned}?
            \end{itemize}
        \pause \item When you're done, try implementing this function recursively (without helper functions).
    \end{itemize}
\end{frame}

\end{document}
